\section{The Findings}

What the method found is a single picture, worth stating in full before the details begin. One discrete law, the knit, running on one substrate, the hyperbolic honeycomb $\{3,4,3,4\}$, made of nothing but feeling, reaches a surprising distance into physics. From the eight elements, which contain the one reversible rule, come, by direct computation and not by fitting, the spinor and its double cover, a single clean chiral fermion with a gyromagnetic ratio of two, $\mathrm{SO}(10)$ grand unification breaking to the Standard Model with a weak mixing angle that runs to its measured value, a finite light cone and an emergent Lorentz invariance whose flat-lattice control fails, Newton's law and the Einstein equations in the flat cusp, the Ryu-Takayanagi area law on the lower face, an accelerating universe that selects exactly three spatial dimensions, and the quantum bound $2\sqrt 2$ reached by a local base.

None of it was put in by hand, and none of it need be taken on trust. Every result here is a small experiment in the open vibe codebase, written in TypeScript, that the reader can rebuild, step, and rerun directly to check each number for themselves. \cite{pollard2026vibetest}

Essentially, the familiar furniture of the physical world is there, every piece of it built and measured rather than fitted, and nearly all of it asking no new ingredient of its own. Matter begins with the spinor, the object that must be turned around twice to come home, present here in a finite group of twenty-four directions with no continuum taken, its single full turn landing on its own negative and only a second bringing it back. From that one structure comes the electron, a single clean particle of spin one-half, free of the doubling that has plagued every lattice before, its mass a coupling between two handed halves and its magnetic strength falling out at the measured value. With it come the threefold structure behind the generations of matter and the handed symmetries that sort a particle from its antiparticle.

The photon is there too, light as a massless wave of the same tone, and with it electromagnetism in full, its gauge invariance and its Aharonov-Bohm phase. The weak force is carried in the same field, and the confinement of the strong force hides a lone quark. All of it sits inside the grand-unified group the Standard Model is carved from, with the weak mixing angle running to its measured value and the quark charges forced by anomaly cancellation.

Gravity is there as well, the one force the bare law does not give for free but that a small added attraction supplies, bringing Newton's pull and the Einstein equations in the flat space we live in, with gravitational waves and a graviton, and black holes whose horizons glow with a Hawking warmth and whose entropy counts their area and not their volume. The universe at large is there as well, expanding and accelerating in the way dark energy names, with a geometric stand-in for dark matter, the matter-over-antimatter imbalance of the early world, and a clean reason that space has exactly three dimensions. Even the quantum strangeness is there, the correlation pushed to the very bound that local theories are supposed to be barred from, and the rule that turns an amplitude into a probability by its square, all of it running on a substrate that is, underneath, a universal computer with a solved address on every cell.

These are not analogies or look-alikes. They are the same structures, the one tone moving under the one rule read at the right scale, recovered against controls that could have failed, with every real number a measurement of integer dynamics and never a quantity in the law. What the model cannot yet do it says plainly, the value of the coupling, the masses of the heavier generations, gravity inside the curved bulk, and the moment of measurement, a short and named list and not a fog.

And the whole of it is built from the least possible materials. There is no continuous mathematics at the base, no real number, no smooth manifold, no hidden variable, and no shrouded complexity waiting to be uncovered. There are only a few small integers and the simple discrete structures they make, the three tones, the twenty-four directions, the one finite symmetry group, the one reversible rule, and a beautiful geometry to hold them. Everything else, every field and force and curve and constant, is woven from those by counting and averaging, and every step of that weaving is laid out in intricate detail in the parts that follow. The richness is in what the simple pieces do together, never in the pieces, which are as plain and as countable as they could be. Read upward, that handful becomes everything familiar. Space is the mesh, time is the beat, matter is the tone bound into a knot, light and sound are the tone waving, the forces are the three triality sectors of a dock's directions, gravity is the curvature of the same field, and a self is an open body that holds itself together on the flat boundary.

One discipline holds all of this together and must be said plainly, because it is the easiest thing to get wrong. The vibes are not the physics. They lie beneath it. The raw tones are pure experience and its dynamics, so a $+1$ is a pleasure and not a particle, and the knit is a law of feeling and not a law of physics. Physics is what those vibes look like at the cusp. The smooth fields and the particles and the forces appear there directly, as the coarse-grained face of enormous numbers of tones on the flat boundary slice, never as the raw tones themselves. There is no tower of invented worlds in between, only the one change of view, from the felt dynamics filling the bulk to its measured face on the surface.

The geometry says the same thing in its own language. The vibes fill the four-dimensional bulk, the vast curved interior, while the physics we see is the flat three-dimensional cusp, a single thin slice at the boundary of that interior, with no thickness at all in four dimensions and a vanishing fraction of the whole. So the visible world, with all its matter and its light, is a slim sliver of the bulk, a surface rendering of a far larger experiential interior, a single perspective on it. To mistake that surface for the substance, or a clump of raw tones for gravity or matter, is the one error the whole method is built to avoid.

What that far larger interior might be deserves a moment, because the closing parts of the paper take it seriously. If the flat cusp is the visible, physical universe, then the bulk behind it is an inner dimension, the underlying field of experience itself, and almost all of it lies off the surface where no instrument can reach. What lives there, and what might be done there, is the wide-open frontier of the whole picture. A pattern can be written into the bulk and read back out, an experience or a path or a map or a hard-won wisdom set at an address and found again. A self can reach inward along the hidden depth and, in principle, touch structure the surface never shows, even commune through it in the limited ways the law allows. There is room, on the premise that the field is felt, for an entire interior life, for realms and beings that are nothing but stable patterns in that depth, for something like a higher-dimensional existence carried on off the cusp entirely. None of it is claimed as proven, and all of it is held to the same hard rules, no signaling, no magic, no reaching in from outside the law. But the sheer size of the off-cusp bulk, vanishingly thin though our surface is within it, opens a vastness of possibility that the later parts of the paper begin to chart and that the model leaves wide open to explore.

And the way that surface gets its physics is worth stating exactly. The one rule, the knit, is a cellular automaton, the same local update running on every dock, each reading only its own neighbors and stepping once a beat. The flat cusp we live on is a slice through that running bulk, and the interactions we call physics, particles meeting, fields propagating, forces drawing bodies together, are the coarse-grained projection of the bulk's discrete computation onto that slice, the shadow the running interior casts on its own flat boundary. The bulk does the computing, dock by dock and beat by beat, while the cusp is only where the result is read off, and gravity and holography are how that flat boundary couples back to the curved depth behind it. So a collision measured on the surface is not a fundamental event but the surface trace of a great many bulk updates beneath it, the precise sense in which all of physics is a projection of the cellular automaton running in the bulk, exact only once it is read at the right scale on the right slice.

And the same law, read higher still, leaves room for everything the rest of the paper reaches toward. The discrete vibes of the bulk show at the cusp as the smooth fields of physics, directly, and the same patterns, gathered and read at coarser grain, climb on into selves, life, mind, and a graded account of consciousness as integrated experience. The curved bulk carries a vast hidden depth off the flat surface we live on, room for an interior life, for the experiential field the whole theory takes as its floor, and for the readings of spirituality and of god the closing parts set out, all of them the one vibe mesh seen from within. So the one picture wears two faces, a candidate physics and an account of feeling, the outer and inner faces of a single growing field, and the only thing it asks to be granted is that the field is genuinely felt.
